\textbf{Result.} Span efficiency at the last station ranges from $0.9673$ on the rigid baseline B to $0.9859$ on F, every value below the planar-wake limit. Induced drag at that station spans $54.93$ to $55.98$ counts.

\textbf{Decay between the first and last plane} is $9.5$ to $10.2$ per cent, and is consistent across the geometries: B $9.46$\,\%, W $9.68$\,\%, F $9.95$\,\%, H $9.98$\,\%, C $10.19$\,\%, M $10.22$\,\%. That consistency matters more than the magnitude, because a decay rate that differed between geometries would mean the comparison was being made on wakes resolved to different degrees, and no difference between them could be attributed to the wing.

\textbf{Against the baseline.} Every candidate measured here carries a higher span efficiency than the rigid baseline's $0.9673$ (C, F, H, M and W). The margin is small, one to two per cent, which is the size of spanload improvement the low-order design study predicted for this class of change.
